% eksperymentalne tłumaczenie na włoski

%

\section{Twierdzenie Pappusa} % lang-pl
\section{Teorema di Pappo} % lang-it

\begin{proposition}
[twierdzenie Pappusa] % lang-pl
[teorema di Pappo] % lang-it
\index{twierdzenie!Pappusa}% % lang-pl
	Dane są dwie różne proste, $k$ oraz $l$ i sześć punktów różnych od $k \cdot l$: $K_1$, $K_2$, $K_3$ na prostej $k$, $L_1$, $L_2$, $L_3$ na prostej $l$. % lang-pl
	Wtedy punkty $K_1L_1 \cdot K_3L_2$, $L_1K_2 \cdot K_3L_3$ oraz $K_2L_2 \cdot K_1L_3$ są współliniowe. % lang-pl
\index{teorema!di Pappo}% % lang-it
	Siano date due rette distinte $k$ e $l$ e sei punti distinti da $k \cdot l$: $K_1$, $K_2$, $K_3$ sulla retta $k$, $L_1$, $L_2$, $L_3$ sulla retta $l$. % lang-it
	Allora i punti $K_1L_1 \cdot K_3L_2$, $L_1K_2 \cdot K_3L_3$ e $K_2L_2 \cdot K_1L_3$ sono allineati. % lang-it
\end{proposition}

% Eves s. 250 

% TODO: https://en.wikipedia.org/wiki/Pappus%27s_hexagon_theorem

Piszą o nim Audin \cite[s. 25, 151, 171]{audin_2003} (w wersji afinicznej, potem rzutowej), Neugebauer \cite[s. 115, 116]{neugebauer_2018}, Hartshorne \cite[s. 131, 132]{hartshorne2000} (wersja afinicza). % lang-pl
Ne parlano Audin \cite[pp. 25, 151, 171]{audin_2003} (prima nella versione affine, poi in quella proiettiva), Neugebauer \cite[pp. 115, 116]{neugebauer_2018} e Hartshorne \cite[pp. 131, 132]{hartshorne2000} (versione affine). % lang-it
% Hartshorne 62?
Hartshorne wspomni, że Hilbert pokazał w 1971, że niekoniecznie przemienne ciało\footnote{Czyli pierścień z dzieleniem.} $F$ jest przemienne wtedy i tylko wtedy, gdy twierdzenie Pappusa zachodzi w płaszczyźnie $F \times F$. % lang-pl
Hartshorne osserva che Hilbert mostrò nel 1971 che un corpo non necessariamente commutativo\footnote{Cioè un anello con divisione.} $F$ è commutativo se e solo se il teorema di Pappo vale nel piano $F \times F$. % lang-it

Warto zwrócić uwagę na podobieństwa między twierdzeniami Pascala i Pappusa. % lang-pl
Vale la pena osservare le somiglianze tra i teoremi di Pascal e di Pappo. % lang-it

% \item Zna przykłady przekształceń rzutowych i umie je stosować w zadaniach i dowodach twierdzeń rzutowych (Desarguesa, Pappusa, Pascala, Brianchona). zna pojęcia: biegun i biegunowa i potrafi formułować twierdzenia dualne.  
% wg Neugebauera, Brianchon jest dualny do Pascala

% Pascal, Brianchon, Chasles -> Eves s. 256 

https://www.deltami.edu.pl/2015/03/dowod-w-stylu-greckim/

%